% Configured surface, achieved tool orientation, and physical surface.
\begin{tikzpicture}[
    scale=1.15,
    reference/.style={red!75!black, line width=1.1pt},
    physical/.style={blue!70!black, line width=1.1pt},
    tool/.style={green!55!black, line width=1.1pt},
    datum/.style={black, line width=0.8pt, densely dashed},
    conceptangle/.style={line width=0.8pt,
      {Latex[length=1.15mm,width=1.10mm]}-{Latex[length=1.15mm,width=1.10mm]}},
    lbl/.style={font=\footnotesize, align=center},
  ]

% Shared legend below the geometry: surfaces first, then tool directions.
% Keep both rows clear of the physical surface and preserve the line styles.
\draw[reference] (-1.30,-1.22) -- (-0.55,-1.22);
\node[lbl, red!75!black, anchor=west] at (-0.40,-1.22)
  {Configured surface};
\draw[physical] (3.25,-1.22) -- (4.00,-1.22);
\node[lbl, blue!70!black, anchor=west] at (4.15,-1.22)
  {Physical surface};
\draw[datum] (-1.30,-1.66) -- (-0.55,-1.66);
\node[lbl, black, anchor=west] at (-0.40,-1.66)
  {Desired orientation};
\draw[tool] (3.25,-1.66) -- (4.00,-1.66);
\node[lbl, green!55!black, anchor=west] at (4.15,-1.66)
  {Tool face};

% The configured surface gives the desired parallel tool direction, carried to
% the tool as a short reference ray. The orientation reached before approach
% retains a small schematic difference from it.
\coordinate (origin) at (0,0);
\coordinate (toolpivot) at (0,1.30);
\coordinate (labelcolumn) at (3.60,0);
% Match the left and right limits of the desired-orientation datum.
\draw[reference] (-2.40,0) -- (2.40,0);
\draw[datum] (-2.40,1.30) -- (2.40,1.30);
\draw[tool, rotate around={10:(0,1.30)}] (-2.10,1.30) -- (2.10,1.30);
% The tool tip, drawn perpendicular to the face where the datum crosses it,
% so the green object reads as the end of a tool and not as a bare line. It
% rotates with the face, so it stays perpendicular if the tilt is retuned.
\draw[tool, rotate around={10:(0,1.30)}] (0,1.30) -- (0,1.65);
\begin{scope}[shift={(toolpivot)}]
  \draw[conceptangle]
    (0:1.90) arc[start angle=0,end angle=10,radius=1.90];
  \coordinate (toolanglemid) at (5:1.90);
\end{scope}
\node[lbl, anchor=west] (toollabel) at (labelcolumn |- toolanglemid)
  {Desired--achieved difference};

% Placement and calibration can leave the physical surface angularly offset.
% Keep the 19-degree slope with the same horizontal endpoints as the datum.
\draw[physical] (-2.40,{-2.40*tan(19)}) -- (2.40,{2.40*tan(19)});
\draw[conceptangle]
  (0:1.90) arc[start angle=0,end angle=19,radius=1.90];
\coordinate (surfaceanglemid) at (9.5:1.90);
\node[lbl, anchor=west] (surfacelabel) at (labelcolumn |- surfaceanglemid)
  {Configured--physical difference};

% Retain the original canvas size, with room for the two-row legend below.
\pgfresetboundingbox
\begin{scope}[reset cm]
  \path[use as bounding box] (-110.1638pt,-60pt)
    rectangle (258.45518pt,56.55481pt);
\end{scope}
\end{tikzpicture}
